\documentclass[11pt,a4paper]{article}

\usepackage[utf8]{inputenc}
\usepackage[T1]{fontenc}
\usepackage[french]{babel}
\usepackage[left=2.2cm,right=2.2cm,top=2.4cm,bottom=2.2cm]{geometry}
\usepackage[default,scale=0.96]{sourcesanspro}
\usepackage[varqu,varl]{inconsolata}
\usepackage{microtype}
\usepackage{xcolor}
\usepackage{listings}
\usepackage{booktabs}
\usepackage{tabularx}
\usepackage{enumitem}
\usepackage{titlesec}
\usepackage{fancyhdr}
\usepackage{tcolorbox}
\usepackage{hyperref}

\tcbuselibrary{skins,breakable}

% ---------------------------------------------------------------- %
% Palette : l'échelle réglementaire du DPE, la même que les visuels
% ---------------------------------------------------------------- %
\definecolor{encre}{HTML}{11161C}
\definecolor{encre2}{HTML}{49555F}
\definecolor{encre3}{HTML}{79838F}
\definecolor{trait}{HTML}{D6DAE1}
\definecolor{fondcode}{HTML}{F5F7F9}
\definecolor{vertA}{HTML}{008550}
\definecolor{orangeF}{HTML}{EF8023}
\definecolor{rougeG}{HTML}{C41520}
\definecolor{bleu}{HTML}{1B5E8C}

\hypersetup{
  colorlinks=true, urlcolor=bleu, linkcolor=encre, citecolor=bleu,
  pdftitle={DPE Lakehouse - Tutoriel d'exploitation},
  pdfauthor={Eben Ezer NGBLOGNI}
}

\titleformat{\section}{\Large\bfseries\color{encre}}{\thesection.}{0.5em}{}[\vspace{-0.6em}\color{trait}\rule{\textwidth}{0.8pt}]
\titleformat{\subsection}{\large\bfseries\color{encre}}{\thesubsection}{0.5em}{}
\titleformat{\subsubsection}{\normalsize\bfseries\color{encre2}}{}{0em}{}
\titlespacing{\section}{0pt}{1.6em}{0.8em}
\titlespacing{\subsection}{0pt}{1.2em}{0.4em}

\pagestyle{fancy}
\fancyhf{}
\fancyhead[L]{\small\color{encre3}DPE Lakehouse}
\fancyhead[R]{\small\color{encre3}\leftmark}
\fancyfoot[C]{\small\color{encre3}\thepage}
\renewcommand{\headrulewidth}{0.4pt}
\renewcommand{\headrule}{\color{trait}\hrule width\headwidth height\headrulewidth}

\setlist[itemize]{leftmargin=1.2em,itemsep=0.25em,topsep=0.35em}
\setlist[enumerate]{leftmargin=1.5em,itemsep=0.3em,topsep=0.35em}

% ---------------------------------------------------------------- %
% Blocs de code
% ---------------------------------------------------------------- %
\lstdefinestyle{base}{
  backgroundcolor=\color{fondcode},
  basicstyle=\ttfamily\footnotesize\color{encre},
  keywordstyle=\color{bleu}\bfseries,
  commentstyle=\color{encre3}\itshape,
  stringstyle=\color{vertA},
  numberstyle=\tiny\color{encre3},
  frame=leftline, framerule=2pt, rulecolor=\color{trait},
  framesep=8pt, xleftmargin=12pt,
  breaklines=true, breakatwhitespace=false,
  showstringspaces=false, tabsize=2,
  columns=fullflexible, keepspaces=true,
  literate=
    {é}{{\'e}}1 {è}{{\`e}}1 {ê}{{\^e}}1 {à}{{\`a}}1 {ç}{{\c c}}1
    {É}{{\'E}}1 {ô}{{\^o}}1 {î}{{\^\i}}1 {û}{{\^u}}1 {ù}{{\`u}}1
    {²}{{\textsuperscript{2}}}1 {€}{{\euro}}1 {—}{{---}}1 {·}{{\textperiodcentered}}1
}
\lstdefinestyle{sh}{style=base,language=bash,
  keywords={make,docker,python,cd,export,git,psql,dbt,curl,sbt},
  morekeywords={compose,up,down,run,exec,ps}}
\lstdefinestyle{py}{style=base,language=Python}
\lstdefinestyle{sc}{style=base,language=Scala}
\lstdefinestyle{sql}{style=base,language=SQL,morekeywords={with,as,filter,over,partition,round,nullif,coalesce}}
\lstdefinestyle{yml}{style=base,language=Python,keywords={}}

\newtcolorbox{pourquoi}[1][]{
  enhanced, breakable, colback=white, colframe=vertA,
  boxrule=0pt, leftrule=2.5pt, arc=0pt, outer arc=0pt,
  left=10pt, right=8pt, top=7pt, bottom=7pt, #1}

\newtcolorbox{attention}[1][]{
  enhanced, breakable, colback=white, colframe=orangeF,
  boxrule=0pt, leftrule=2.5pt, arc=0pt, outer arc=0pt,
  left=10pt, right=8pt, top=7pt, bottom=7pt, #1}

\newcommand{\cmd}[1]{\texttt{\small #1}}
\newcommand{\titrebloc}[1]{\textbf{\color{encre}#1}\par\vspace{0.25em}}

\begin{document}

% ================================================================== %
\begin{titlepage}
\thispagestyle{empty}
\vspace*{2.5cm}
{\color{encre3}\small\ttfamily PIPELINE DE DONNÉES · ADEME · 2026}\\[1.2em]
{\Huge\bfseries\color{encre} DPE Lakehouse}\\[0.5em]
{\LARGE\color{encre2} Tutoriel d'exploitation}\\[2em]
{\large\color{encre2} Comment lancer chaque composant, pourquoi il est là,\\ et ce qu'il faut savoir en dire.}\\[3em]

\color{trait}\rule{\textwidth}{0.8pt}\\[1em]
\color{encre2}
\begin{tabularx}{\textwidth}{@{}lX@{}}
\textbf{Source} & ADEME, base des DPE logements existants \\[0.3em]
\textbf{Volumes} & 15,3 M de lignes ingérées, 7 849 191 en entrepôt \\[0.3em]
\textbf{Stack} & Python, Spark/Scala, PostgreSQL, dbt, Airflow, Docker, Power BI \\[0.3em]
\textbf{Dépôt} & \href{https://github.com/ebenezer-ngblogni/dpe-lakehouse}{github.com/ebenezer-ngblogni/dpe-lakehouse} \\
\end{tabularx}\\[1em]
\color{trait}\rule{\textwidth}{0.8pt}

\vfill
{\color{encre}\large\bfseries Eben Ezer NGBLOGNI}\\
{\color{encre3} Élève ingénieur informatique, EILCO}
\end{titlepage}

\tableofcontents
\newpage

% ================================================================== %
\section{Pourquoi ce projet}

\subsection{Le problème choisi}

La France compte environ 5,8 millions de logements classés F ou G, les
\textbf{passoires thermiques}. La loi Climat et Résilience les interdit
progressivement à la location. Savoir \emph{où} elles se trouvent et
\emph{à quoi} elles ressemblent est une question publique, et l'ADEME publie
la matière première en accès libre.

\begin{pourquoi}
\titrebloc{Ce que je réponds si on me demande « pourquoi ce sujet ? »}
Trois raisons, dans cet ordre :

\begin{enumerate}
  \item \textbf{Le volume est réel.} 15,3 millions de lignes obligent à traiter
        en distribué. Sur 50 000 lignes, Spark serait une pose ; ici il est
        nécessaire.
  \item \textbf{La donnée est sale, et de façon intéressante.} Codes postaux
        à la place des codes INSEE, communes fusionnées préfixées
        \cmd{old}, noms de colonnes corrompus par l'export CSV. Ce sont des
        problèmes qu'on ne rencontre pas dans un jeu de données pédagogique.
  \item \textbf{Le résultat se raconte.} La part de passoires s'effondre après
        la réglementation thermique de 1974. C'est un fait mesuré sur des
        données de terrain, pas une illustration.
\end{enumerate}
\end{pourquoi}

\subsection{Ce que le projet démontre}

\begin{center}
\begin{tabularx}{\textwidth}{@{}lX@{}}
\toprule
\textbf{Compétence} & \textbf{Où elle se voit} \\
\midrule
Ingestion résiliente & Pagination par curseur, partitions mensuelles, reprise sans doublon \\
Traitement distribué & Spark 3.5 en Scala, fenêtres, jointures, 15,3 M de lignes \\
Qualité de données & 720 483 rejets tracés avec leur motif, jamais silencieux \\
Modélisation & Schéma en étoile dbt, agrégats commune et département \\
Orchestration & DAG Airflow à six tâches, chacune reprenable \\
Restitution & Rapport Power BI, carte de France à 96 points lisibles \\
Rigueur & 74 tests automatisés, intégration continue GitHub Actions \\
\bottomrule
\end{tabularx}
\end{center}

\subsection{L'architecture en une phrase}

Une \textbf{architecture médaillon} : la donnée brute atterrit en
\emph{bronze} sans transformation, est nettoyée en \emph{silver}, puis
modélisée en \emph{gold} (les marts) pour la restitution.

\begin{pourquoi}
\titrebloc{Pourquoi trois couches et pas une seule}
Parce que chaque couche répond à une question différente.

\textbf{Bronze} répond à « qu'a dit la source, exactement ? ». On n'y touche
à rien : si une règle métier change, on rejoue depuis le bronze sans
retélécharger 15 millions de lignes — l'API met environ cinq heures.

\textbf{Silver} répond à « quelles lignes sont exploitables ? ». C'est là que
vivent la déduplication et les règles qualité.

\textbf{Gold} répond à « que veut voir l'analyste ? ». Les agrégats y sont
pré-calculés pour que Power BI réponde instantanément.
\end{pourquoi}

\newpage
% ================================================================== %
\section{Démarrer la pile}

\subsection{Prérequis}

\begin{lstlisting}[style=sh]
# Depuis la racine du dépôt
cd ~/Bureau/lab/DataEng

# Le lac vit sur une partition séparée (1 To). Si `data` est un lien mort :
udisksctl mount -b /dev/sda5
ls data/          # doit lister bronze/ silver/ exports/

# Environnement Python
make setup        # crée .venv et installe les dépendances
\end{lstlisting}

\begin{attention}
\titrebloc{Le piège des identifiants Docker}
Les conteneurs écrivent dans des volumes montés depuis l'hôte. Si l'UID du
conteneur ne correspond pas au vôtre, l'écriture échoue avec un
\cmd{Permission denied} peu bavard. Il faut donc exporter ces variables
\emph{avant} de lancer la pile :

\begin{lstlisting}[style=sh]
export AIRFLOW_UID=$(id -u)
export HOST_UID=$(id -u)
export DOCKER_GID=$(getent group docker | cut -d: -f3)
\end{lstlisting}

\cmd{DOCKER\_GID} est nécessaire parce qu'Airflow soumet les jobs Spark via
\cmd{docker exec} : il lui faut le droit de parler au démon Docker.
\end{attention}

\subsection{Lancer les services}

\begin{lstlisting}[style=sh]
make up           # équivalent de : docker compose up -d
docker compose ps # vérifier que tout est "healthy"
\end{lstlisting}

\begin{center}
\begin{tabularx}{\textwidth}{@{}llX@{}}
\toprule
\textbf{Service} & \textbf{Adresse} & \textbf{Rôle} \\
\midrule
Airflow    & \url{http://localhost:8080} & Orchestration (admin / admin) \\
MinIO      & \url{http://localhost:9003} & Stockage objet compatible S3 \\
PostgreSQL & \cmd{localhost:5434}        & Entrepôt analytique \\
dbt docs   & \url{http://localhost:8081} & Documentation et lignage \\
\bottomrule
\end{tabularx}
\end{center}

\begin{attention}
\titrebloc{Pourquoi ces ports décalés}
5434 et non 5432, 9002/9003 et non 9000/9001 : les ports par défaut étaient
déjà pris par un PostgreSQL local et un autre projet. Le décalage est
volontaire et propagé partout — \cmd{docker-compose.yml},
\cmd{profiles.yml}, le \cmd{Makefile}, le code Scala. C'est typiquement le
genre de détail qu'un recruteur apprécie de vous entendre expliquer.
\end{attention}

\newpage
% ================================================================== %
\section{Étape 1 — Ingestion (Python)}

\subsection{Lancer}

\begin{lstlisting}[style=sh]
make ingest-sample   # un seul mois, environ 2 minutes — pour tester
make ingest-full     # tout l'historique depuis juillet 2021, environ 5 h
make ingest-status   # compare l'état local à la source, partition par partition
\end{lstlisting}

\subsection{Pourquoi Python ici, et pas Spark}

\begin{pourquoi}
L'ingestion est un travail \textbf{d'entrées-sorties réseau}, pas de calcul.
On attend l'API, on écrit sur disque, on recommence. Spark n'apporterait rien :
son intérêt est de paralléliser du calcul sur des données déjà présentes.
Python avec PyArrow suffit et se déploie en trois lignes.
\end{pourquoi}

\subsection{Les deux idées à retenir}

\subsubsection{Pagination par curseur}

L'API ADEME (Data Fair) plafonne à 10 000 lignes par requête et fournit le
lien suivant dans un en-tête HTTP \cmd{Link}. On ne calcule donc jamais
d'\emph{offset} soi-même : on suit le curseur.

\begin{lstlisting}[style=py]
# Le lien `next` embarque déjà tous les paramètres de la requête.
next_link = response.links.get("next")
url = next_link.get("url") if next_link else None
\end{lstlisting}

\begin{pourquoi}
\titrebloc{Pourquoi c'est mieux qu'un offset}
Avec \cmd{OFFSET n}, si une ligne est insérée en cours de parcours, tout
décale et on saute ou duplique des enregistrements. Le curseur est stable :
il pointe une position logique dans le jeu, pas un rang.
\end{pourquoi}

\subsubsection{Promotion atomique}

Chaque mois est écrit dans un répertoire de \emph{staging}, puis promu d'un
seul coup, accompagné d'un manifeste :

\begin{lstlisting}[style=py]
def is_partition_complete(self, table, year, month) -> bool:
    """Une partition est réputée chargée si son manifeste est présent
    ET complet."""
    manifest = self.read_manifest(table, year, month)
    return manifest is not None and manifest.is_complete
\end{lstlisting}

\begin{pourquoi}
\titrebloc{Le problème que ça résout}
Si le processus meurt à 80 \% d'un mois, une écriture directe laisserait une
partition à moitié remplie — indiscernable d'une partition complète. À la
reprise, le pipeline la croirait faite et perdrait 20 \% des données,
silencieusement.

Avec la promotion atomique, une partition incomplète n'a pas de manifeste :
elle est simplement rejouée. \textbf{C'est ce qui rend l'ingestion
idempotente} : la relancer dix fois donne le même résultat qu'une fois.
\end{pourquoi}

\subsubsection{Le choix du format de transport}

\begin{pourquoi}
\titrebloc{Une décision prise après mesure, pas par principe}
L'API sert du JSON ou du CSV. J'ai mesuré les deux : le CSV divise
\textbf{par 2,8} le temps de chargement, parce qu'il évite de construire un
objet Python par ligne.

Le prix à payer : l'export CSV de l'API corrompt certains noms de colonnes en
remplaçant un underscore par une espace — le schéma déclare
\cmd{conso\_5\_usages\_par\_m2\_ef}, le CSV émet
\cmd{conso\_5 usages\_par\_m2\_ef}. On répare à la lecture, et trois tests de
non-régression verrouillent la correction.
\end{pourquoi}

\newpage
% ================================================================== %
\section{Étape 2 — Nettoyage (Spark / Scala)}

\subsection{Lancer}

\begin{lstlisting}[style=sh]
make build        # compile le jar avec sbt-assembly
make test-scala   # 15 tests unitaires
make silver       # bronze -> silver, sur la totalité du lac
\end{lstlisting}

\subsection{Pourquoi Scala et pas PySpark}

\begin{pourquoi}
Trois arguments, du plus fort au plus faible :

\begin{enumerate}
  \item \textbf{Scala est la langue native de Spark.} Le moteur est écrit en
        Scala ; on parle directement à l'API, sans passerelle Py4J.
  \item \textbf{Le typage attrape les erreurs à la compilation.} Une faute de
        frappe sur un nom de colonne ou un mauvais type est signalée avant de
        lancer un job qui tourne vingt minutes.
  \item \textbf{C'est ce que demandent les offres.} Beaucoup d'équipes data
        historiques sont en Scala ; savoir y écrire ouvre des portes que
        PySpark seul n'ouvre pas.
\end{enumerate}

L'argument que je n'utilise \emph{pas} : « c'est plus rapide ». Depuis les
DataFrames, PySpark et Scala passent par le même optimiseur Catalyst et les
performances sont comparables. Le dire serait faux, et un bon recruteur le
sait.
\end{pourquoi}

\subsection{Les trois transformations}

\subsubsection{Déduplication par fenêtre}

\begin{lstlisting}[style=sc]
val fenetre = Window
  .partitionBy(col("numero_dpe"))
  .orderBy(
    col("date_derniere_modification_dpe").desc_nulls_last,
    col("date_reception_dpe").desc_nulls_last
  )

df.withColumn("_rang", row_number().over(fenetre))
  .filter(col("_rang") === 1)
  .drop("_rang")
\end{lstlisting}

\begin{attention}
\titrebloc{Résultat mesuré : 0 ligne supprimée}
J'attendais des milliers de doublons. Il n'y en a aucun — le numéro de DPE est
réellement unique dans le lot ingéré. \textbf{J'ai gardé le code et documenté
le résultat} plutôt que de le retirer : c'est un garde-fou qui protège une
hypothèse, et savoir qu'elle tient a de la valeur.

Si on me demande « pourquoi garder du code qui ne fait rien ? », la réponse
est : il ne fait rien \emph{aujourd'hui}. Le jour où l'ADEME republie un
diagnostic, il fera son travail.
\end{attention}

\subsubsection{Règles qualité avec motif conservé}

\begin{lstlisting}[style=sc]
private def motifRejet: Column =
  when(col("numero_dpe").isNull, "identifiant_absent")
    .when(!col("etiquette_dpe").isin(EtiquettesValides: _*), "etiquette_dpe_invalide")
    .when(col("date_etablissement_dpe").isNull, "date_etablissement_absente")
    .when(col("code_insee_ban").isNull, "commune_absente")
    .when(col("surface_habitable_logement") <= 0, "surface_invalide")
    // Un logement de plus de 1 000 m2 releve de la saisie erronee.
    .when(col("surface_habitable_logement") > 1000, "surface_aberrante")
    .when(col("conso_5_usages_par_m2_ep") < 0, "consommation_invalide")
\end{lstlisting}

\begin{pourquoi}
\titrebloc{Pourquoi écrire les rejets au lieu de les jeter}
\textbf{720 483 lignes} sont écartées. Si on se contentait d'un
\cmd{filter}, elles disparaîtraient sans trace et personne ne saurait
pourquoi le total ne tombe pas juste.

Ici chaque rejet garde son motif et part dans \cmd{silver/dpe\_rejets}. On
peut donc répondre à « combien de lignes perd-on, et pour quelle raison ? » —
question que pose systématiquement un responsable data.
\end{pourquoi}

\subsubsection{Le seuil de géocodage, corrigé après mesure}

\begin{lstlisting}[style=sc]
val SeuilScoreBan: Double = 0.64  // mediane mesuree (etait 0.8 : 79 % marques)
\end{lstlisting}

\begin{attention}
J'avais fixé le seuil à 0{,}8 « parce que ça semblait raisonnable ». Résultat :
79 \% des lignes signalées comme douteuses — un filtre qui marque quatre lignes
sur cinq ne filtre rien. J'ai regardé la distribution réelle, pris la médiane,
et surtout basculé sur le \emph{statut} de géocodage fourni par la source,
qui est un verdict et non un score.
\end{attention}

\newpage
% ================================================================== %
\section{Étape 3 — Entrepôt (PostgreSQL)}

\subsection{Lancer}

\begin{lstlisting}[style=sh]
make warehouse    # silver -> PostgreSQL, fenêtre depuis 2024

# Se connecter pour vérifier
docker exec -it dpe_warehouse psql -U dpe -d dpe
\end{lstlisting}

\begin{lstlisting}[style=sql]
-- Quelques requêtes de contrôle
SELECT count(*) FROM silver.dpe_courant;        -- 7 849 191
SELECT count(*) FROM silver.dpe_rejets;         --   720 483
\d marts.*
\end{lstlisting}

\subsection{Pourquoi PostgreSQL et pas Snowflake ou BigQuery}

\begin{pourquoi}
\textbf{Parce que le volume ne le justifie pas.} 7,8 millions de lignes sur
40 colonnes tiennent dans 2,7 Go. Un PostgreSQL sur un poste répond en
quelques secondes. Payer un entrepôt cloud pour ça serait une dépense sans
contrepartie.

Ce que je dis ensuite, parce que c'est la vraie question : \textbf{à partir de
quand je changerais}. Trois signaux — quand plusieurs équipes requêtent en
concurrence, quand les volumes dépassent ce qu'une machine tient en mémoire,
ou quand il faut séparer stockage et calcul pour maîtriser les coûts. Aucun
n'est atteint ici.
\end{pourquoi}

\subsection{Deux optimisations concrètes}

\begin{itemize}
  \item \textbf{Projection de colonnes} : 40 colonnes chargées sur 78
        disponibles. L'entrepôt passe de 5 825 Mo à 2 700 Mo. On ne charge que
        ce que les marts consomment.
  \item \textbf{Fenêtre temporelle} : seules les données depuis 2024 partent
        en entrepôt. Le lac, lui, garde tout l'historique. C'est un arbitrage
        coût de stockage contre profondeur d'analyse, assumé et documenté.
\end{itemize}

\newpage
% ================================================================== %
\section{Étape 4 — Modélisation (dbt)}

\subsection{Lancer}

\begin{lstlisting}[style=sh]
cd dbt/dpe_analytics

dbt deps                        # installe dbt_utils
dbt seed --profiles-dir .       # charge le référentiel des départements
dbt run  --profiles-dir .       # construit les modèles
dbt test --profiles-dir .       # 44 tests de contrat
dbt docs generate --profiles-dir .
dbt docs serve --profiles-dir . --port 8081
\end{lstlisting}

\subsection{L'interface : le graphe de lignage}

\cmd{dbt docs serve} ouvre une interface web sur
\url{http://localhost:8081}. Deux choses à savoir montrer :

\begin{enumerate}
  \item \textbf{Le graphe de lignage} (bouton en bas à droite). Il affiche
        automatiquement les dépendances entre modèles :
        \cmd{stg\_dpe} $\rightarrow$ \cmd{fct\_dpe} et \cmd{dim\_commune}
        $\rightarrow$ les trois marts. \textbf{Ce graphe n'est pas dessiné à la
        main} : dbt le déduit des \cmd{ref()} dans le SQL.
  \item \textbf{La fiche d'un modèle} : sa description, ses colonnes, ses
        tests, et le SQL compilé. C'est la documentation qui ne peut pas
        mentir, puisqu'elle est générée depuis le code.
\end{enumerate}

\subsection{Pourquoi dbt}

\begin{pourquoi}
\titrebloc{Ce que dbt apporte qu'un script SQL n'apporte pas}
\begin{itemize}
  \item \textbf{Le lignage gratuit.} On écrit \cmd{\{\{ ref('fct\_dpe') \}\}}
        au lieu du nom de table ; dbt en déduit l'ordre d'exécution et le
        graphe de dépendances.
  \item \textbf{Les tests comme contrat.} Unicité, non-nullité, plages de
        valeurs, cohérence arithmétique — déclarés en YAML à côté du modèle,
        pas dans un script séparé qu'on oublie de lancer.
  \item \textbf{La documentation versionnée.} Elle vit dans le dépôt et suit
        le code.
\end{itemize}
\end{pourquoi}

\subsection{Le piège de la moyenne de taux}

\begin{lstlisting}[style=sql]
-- CORRECT : rapport de sommes
round(100.0 * sum(nb_passoires) / nullif(sum(nb_dpe), 0), 2) AS taux_passoires_pct

-- FAUX : moyenne de taux, dès que les dénominateurs diffèrent
avg(taux_passoires_pct)
\end{lstlisting}

\begin{attention}
\titrebloc{À savoir expliquer par cœur}
Une commune de 30 diagnostics et une préfecture de 40 000 ne pèsent pas
pareil. Faire la moyenne de leurs taux donnerait autant de poids à l'une
qu'à l'autre.

C'est pour cette raison que \textbf{les taux sont des colonnes pré-calculées
en SQL} et non des mesures Power BI : le modèle sémantique est en lecture
seule, donc je ne peux pas y écrire de DAX. J'ai fait remonter le calcul dans
dbt, où il est juste et testé.
\end{attention}

\subsection{Les deux tests rouges, assumés}

\begin{pourquoi}
Sur 44 tests dbt, 42 passent et 2 échouent. \textbf{Je les laisse en rouge.}

Ils signalent un vrai doublon dans la source : l'ADEME a publié deux fois le
même diagnostic, une version incomplète puis une version complète. Le test dit
la vérité ; l'assouplir pour obtenir un tableau vert reviendrait à masquer un
défaut réel.

Si on me demande « pourquoi ne pas corriger ? » : parce que je n'ai pas encore
tranché la cause racine, et qu'un correctif posé sans comprendre en
créerait un autre. En attendant, le test documente le problème.
\end{pourquoi}

\newpage
% ================================================================== %
\section{Étape 5 — Orchestration (Airflow)}

\subsection{Lancer}

\begin{lstlisting}[style=sh]
make up   # démarre webserver + scheduler
# Interface : http://localhost:8080   (admin / admin)
\end{lstlisting}

\subsection{L'interface : ce qu'il faut savoir montrer}

\begin{center}
\begin{tabularx}{\textwidth}{@{}lX@{}}
\toprule
\textbf{Vue} & \textbf{Ce qu'elle apprend} \\
\midrule
\textbf{Grid}  & Historique des exécutions. Chaque carré est une tâche : vert réussi, rouge échoué. On lit d'un coup d'{\oe}il quelle étape casse et depuis quand. \\
\textbf{Graph} & Le DAG lui-même, avec les dépendances. C'est la vue à ouvrir en entretien. \\
\textbf{Logs}  & Sortie de chaque tâche. C'est là qu'on débogue. \\
\textbf{Code}  & Le fichier Python du DAG, tel que le scheduler le lit. \\
\bottomrule
\end{tabularx}
\end{center}

\subsection{Le DAG, en six lignes}

\begin{lstlisting}[style=py]
debut >> ingestion >> transformation >> chargement \
      >> dbt_build >> dbt_test >> dbt_docs >> fin
\end{lstlisting}

Chaque tâche est un \cmd{BashOperator} qui soumet un travail dans le
conteneur adéquat :

\begin{lstlisting}[style=py]
transformation = BashOperator(
    task_id="spark_bronze_vers_silver",
    bash_command=(
        "docker exec dpe_spark spark-submit "
        "--class fr.dpelab.silver.DpeSilverJob --master 'local[*]' "
        f"--conf 'spark.local.dir={SPARK_DATA_ROOT}/tmp-spark' "
        f"{SPARK_JAR} --bronze-path {SPARK_BRONZE}/dpe_existant ..."
    ),
)
\end{lstlisting}

\begin{pourquoi}
\titrebloc{Pourquoi Airflow, et pas un cron}
Un cron sait lancer une commande à une heure. Il ne sait pas :

\begin{itemize}
  \item \textbf{gérer les dépendances} — ne pas charger l'entrepôt si le
        nettoyage a échoué ;
  \item \textbf{rejouer une seule étape} sans refaire les précédentes ;
  \item \textbf{montrer l'historique} : quelle exécution a échoué, quand,
        avec quels logs.
\end{itemize}

C'est exactement ce qu'apporte Airflow, et c'est pourquoi il est devenu le
standard de fait.
\end{pourquoi}

\begin{attention}
\titrebloc{Pourquoi \cmd{docker exec} plutôt qu'un opérateur Spark}
Airflow tourne dans son conteneur, Spark dans le sien. Plutôt que d'installer
Spark dans l'image Airflow — ce qui la ferait grossir et mêlerait deux
responsabilités — Airflow \emph{soumet} le travail au conteneur Spark. Chaque
image reste spécialisée. En production on utiliserait le
\cmd{SparkSubmitOperator} ou un cluster Kubernetes ; le principe est le même.
\end{attention}

\subsection{La fenêtre glissante}

Le DAG ne rejoue pas les 15 millions de lignes chaque nuit : il traite une
fenêtre glissante de trois mois. L'ADEME corrigeant parfois des diagnostics
publiés, on repasse sur le passé récent sans tout refaire.

\newpage
% ================================================================== %
\section{Étape 6 — Restitution (Power BI)}

\subsection{Lancer}

\begin{lstlisting}[style=sh]
make powerbi-export   # exporte les marts en CSV

# ou de façon ciblée
python scripts/export_powerbi.py \
  --table mart_performance_departement --destination powerbi/data

git add powerbi/data && git commit && git push
\end{lstlisting}

\begin{attention}
\titrebloc{Le rapport lit GitHub, pas votre machine}
Faute de licence OneDrive pour téléverser un fichier, Power BI Service lit les
CSV par leur \textbf{URL brute GitHub}. Conséquence directe :
\textbf{l'actualisation dépend d'un \cmd{git push}}. C'est une limite du
montage, pas un oubli — et il vaut mieux la dire soi-même que se la faire
remarquer.
\end{attention}

\subsection{Les deux pièges Power BI, silencieux tous les deux}

\begin{enumerate}
  \item \textbf{Le délimiteur n'est pas détecté.} Tout atterrit dans une
        colonne unique. Il faut le forcer à « Virgule ».
  \item \textbf{L'auto-typage en locale française casse les décimaux.} Le CSV
        écrit \cmd{8.69} avec un point, la locale attend une virgule, et Power
        Query bascule la colonne en texte — latitude et longitude comprises,
        ce qui vide la carte \emph{sans message d'erreur}.
\end{enumerate}

\begin{lstlisting}[style=sh]
Table.TransformColumnTypes(#"En-têtes promus",
  {{"code_departement", type text},
   {"nb_dpe", Int64.Type},
   {"taux_passoires_pct", type number},
   {"latitude", type number}, {"longitude", type number}}, "en-US")
\end{lstlisting}

\begin{pourquoi}
\titrebloc{Et le code département reste du texte}
Typé en nombre, « 01 » perdrait son zéro initial et « 2A » ferait échouer la
conversion. Le même piège existe côté dbt : le seed \cmd{ref\_departement}
déclare ses types explicitement.
\end{pourquoi}

\subsection{Pourquoi la carte est départementale}

\begin{pourquoi}
La première version cartographiait les communes. À 23 628 communes
géolocalisées, chaque point occupe moins d'un pixel : les cercles se
recouvrent, plus rien n'est identifiable.

J'ai essayé le rayon, l'opacité, la carte thermique. Aucun réglage ne corrige
cela — \textbf{c'est le grain qui était en cause}, pas l'affichage. À 96
départements, chaque bulle redevient lisible.

C'est la leçon que je retiens : quand un visuel ne passe pas, il faut souvent
remonter au modèle, pas régler le visuel.
\end{pourquoi}

\newpage
% ================================================================== %
\section{Les résultats à connaître}

\subsection{Les quatre chiffres qui portent le projet}

\begin{center}
\begin{tabularx}{\textwidth}{@{}lX@{}}
\toprule
\textbf{Résultat} & \textbf{Ce qu'il dit} \\
\midrule
\textbf{8,69 \%} & Part de passoires (F ou G) dans le parc diagnostiqué, soit 681 719 logements \\
\textbf{17,5 $\rightarrow$ 0,0 \%} & Part de passoires selon la période de construction, d'avant 1948 à après 2013 \\
\textbf{Facteur 17} & Entre la Creuse (34,6 \%) et l'Hérault (2,1 \%) \\
\textbf{6 € contre 50 €} & Coût annuel de chauffage au m², d'une étiquette A à une étiquette G \\
\bottomrule
\end{tabularx}
\end{center}

\subsection{Le résultat le plus fort}

La chute après \textbf{1974} — année de la première réglementation thermique,
adoptée dans la foulée du choc pétrolier. Avant 1948, 17,5 \% de passoires.
Après 2013, 0,0 \%, et \textbf{60,3 \% des maisons neuves décrochent une
étiquette A ou B}.

C'est l'effet d'une politique publique, mesuré sur des données de terrain.

\subsection{Les limites, à énoncer soi-même}

\begin{attention}
\titrebloc{Biais de sélection}
Un DPE se réalise à la vente ou à la location. Ces 7,8 millions de lignes
décrivent le parc \emph{qui change de mains}, pas le parc français. Les
logements occupés durablement par leurs propriétaires sont sous-représentés —
et ce sont souvent les plus anciens.
\end{attention}

\begin{attention}
\titrebloc{Une corrélation qu'il ne faut pas surinterpréter}
À l'échelle du logement, le type de bâtiment pèse : 14,5 \% de passoires pour
les maisons contre 6,2 \% pour les appartements. Mais à l'échelle
départementale, le lien s'effondre — \textbf{r = 0,31}, soit 9 \% de variance
expliquée. L'agrégation dilue l'effet et le climat reprend le dessus.

Le titre du visuel le dit ainsi, plutôt que de laisser croire à un lien fort.
\end{attention}

\newpage
% ================================================================== %
\section{Questions probables en entretien}

\subsection{« Pourquoi Spark pour 15 millions de lignes ? Pandas suffirait. »}

Non, et c'est mesurable : la déduplication brasse plus de 10 Go de
\emph{shuffle}. Sur un poste à 16 Go de RAM, Pandas sature. Spark décharge sur
disque et termine. J'ai d'ailleurs dû rediriger \cmd{spark.local.dir} vers le
disque du lac, parce que les fichiers temporaires remplissaient le disque
système.

\subsection{« Comment garantissez-vous l'idempotence ? »}

Trois mécanismes, un par couche :

\begin{itemize}
  \item \textbf{Ingestion} : promotion atomique par partition, avec manifeste
        de complétude.
  \item \textbf{Silver} : réécriture complète en
        \cmd{partitionOverwriteMode=dynamic} — seules les partitions produites
        sont remplacées.
  \item \textbf{Entrepôt} : \cmd{SaveMode.Overwrite} avec
        \cmd{truncate=true}, qui vide et recharge sans détruire les index.
\end{itemize}

Rejouer le DAG ne duplique rien.

\subsection{« Qu'est-ce qui vous a le plus surpris ? »}

Que \textbf{trois de mes hypothèses soient fausses}, et de façon mesurable :
la déduplication ne supprime rien, les chaînes de remplacement représentent
0,01 \% et non 8,9 \%, et mon seuil de géocodage marquait 79 \% des lignes.

Je les ai documentées dans le README plutôt que de les effacer. Un pipeline
qui ne raconte que ses succès n'apprend rien à celui qui le reprend.

\subsection{« Que feriez-vous différemment ? »}

Trois choses :

\begin{enumerate}
  \item \textbf{Intégrer l'export Power BI au DAG.} Il est aujourd'hui manuel
        parce qu'il dépend d'un \cmd{git push} ; il faudrait un jeton
        d'écriture dans l'orchestrateur.
  \item \textbf{Trancher le doublon en silver.} Deux tests sont rouges et je
        sais pourquoi, mais pas encore d'où vient le doublon exactement.
  \item \textbf{Passer les marts en incrémental.} Aujourd'hui ils sont
        reconstruits entièrement à chaque exécution ; à volume constant c'est
        acceptable, mais ça ne passera pas l'échelle.
\end{enumerate}

\subsection{« Montrez-moi une décision technique que vous regrettez. »}

Avoir fixé le seuil de géocodage à 0,8 sans regarder la distribution. C'était
une valeur \emph{qui semblait raisonnable}, ce qui n'est pas un critère. La
mesure a montré qu'elle marquait quatre lignes sur cinq. Depuis, je vérifie
un seuil avant de l'écrire.

\vfill
\begin{center}
\color{trait}\rule{0.4\textwidth}{0.8pt}\\[0.8em]
\color{encre3}\small
Document généré depuis le dépôt \cmd{dpe-lakehouse}.\\
Les chiffres sont relevés en base, pas recopiés.
\end{center}

\end{document}
